\begin{proof}[Proof of Theorem~\ref{thm:cr-oracle}]
By Lemma~\ref{lem:cr-endpoint} the conditional minimum is attained at some
$y(z)$, and by Proposition~\ref{prop:cr-cut} minimizing over $z$ is a
minimum cut problem. Every feasible flow has value at most the capacity of
every cut (weak duality), so a flow of value $\operatorname{cap}(S_z)$
proves that $S_z$ is a minimum cut; checking it needs the capacity bounds,
conservation at each residual node and one comparison. A maximum flow and a
minimum cut of equal value exist and are found, for instance, by the
Edmonds--Karp algorithm \cite{EdmondsKarp1972}, whose number of arithmetic
operations is polynomial in the number of nodes and arcs, here
$\abs{\mathcal R}+2$ and at most $2\abs{\mathcal R}$ plus twice the number
of nonzero residual couplings. After multiplication by a common denominator
of polynomial bit length, the capacities are integers, every augmentation is
by an integer, and every flow value is an integer at most the total
capacity, so all intermediate numbers have polynomial bit length. The
coefficients of Proposition~\ref{prop:cr-cut}(b) are obtained by exact
rational evaluation.
\end{proof}

\begin{proof}[Proof of Theorem~\ref{thm:cr-search}]
Let $C$ be a cell of side $h_j$ and $x\in X$ with $x_{\mathcal K}\in C$.
Lemma~\ref{lem:cr-cell} with $B=\mathcal K$ gives
$F(x)\ge\min\{V(v):v\text{ a corner of }C\}-e_{\mathcal K}(C)\ge\beta_V(C)$,
because the core coordinates are continuous and $e_{\mathcal K}(C)\le e_j$.

(a) The unit cube contains $x^*_{\mathcal K}$. If $C\in\mathcal Q_j$ contains
it, then $\beta_V(C)\le F(x^*)=\OPT\le U$, so $C$ is retained and every
child of $C$ that contains $x^*_{\mathcal K}$ belongs to $\mathcal Q_{j+1}$.
(b), (c) These are Theorem~\ref{thm:cr-filter} with $B=\mathcal K$, the
cells $\mathcal Q_j$, $e=e_j$, $\underline V=V$ and $\varepsilon_{\mathrm{or}}=0$; its hypothesis
that some cell of $\mathcal Q_j$ contains $x^*_{\mathcal K}$ holds by (a), and
$\lambda_j\le\OPT$ by the first paragraph applied to $x^*$.
(d) Children of retained cells cover them, so the cells of $\mathcal Q_j$
together with the cells discarded at levels $i<j$ cover $[0,1]^k$. If
$x_{\mathcal K}$ lies in a cell of $\mathcal Q_j$, then $F(x)\ge\lambda_j$ by
the first paragraph. If it lies in a cell $C$ discarded at level $i$, then
$F(x)\ge\beta_V(C)>U^{(i)}\ge U$, where $U^{(i)}$ is the incumbent value at
level $i$. Each step uses only recorded values.
\end{proof}

\begin{proof}[Proof of Proposition~\ref{prop:cr-growth}]
Let $N_j$ be the number of points $v\in h_j\Z^k\cap[0,1]^k$ with
$V(v)\le\OPT+2e_j$. Level $0$ makes $2^k$ queries. By
Theorem~\ref{thm:cr-search}(c), every retained level-$j$ cell has a corner
counted in $N_j$, and a point is a corner of at most $2^k$ cells, so
$\abs{\mathcal Q_j'}\le2^kN_j$. Each retained cell has $2^k$ children with
$2^k$ corners each, so level $j+1$ makes at most $8^kN_j$ queries. A point
counted in $N_j$ satisfies $g\norm{v-v^*}^2\le2e_j=kLh_j^2/4$, so each of
its coordinates lies within $\frac{h_j}2\sqrt{k\kappa_{\mathcal K}}$ of that
of $v^*$; an interval of length $h_j\sqrt{k\kappa_{\mathcal K}}$ contains at
most $1+\sqrt{k\kappa_{\mathcal K}}$ points of $h_j\Z$. Hence
$N_j\le(1+\sqrt{k\kappa_{\mathcal K}})^k$; sum over $j<J$.
\end{proof}
